\documentclass[11pt,a4paper]{article}
\usepackage[utf8]{inputenc}
\usepackage[T1]{fontenc}
\usepackage[french]{babel}
\usepackage{amsmath}
\usepackage{amssymb}
\usepackage{geometry}
\usepackage{hyperref}
\geometry{hmargin=2.5cm,vmargin=2.5cm}

\title{Documentation Technique : Modèle Radiatif à 2 Couches et Dépendance Chimique}
\author{Groupe : [Animal] [Épithète]}
\date{\today}

\begin{document}

\maketitle

\begin{abstract}
Ce document explicite la physique et les mathématiques du script Python modélisant l'équilibre radiatif de la Terre avec une atmosphère discrétisée en deux couches. Contrairement au modèle à une couche, l'émissivité n'est plus arbitraire : elle est calculée dynamiquement via la loi de Beer-Lambert en fonction de la concentration en $CO_2$ et de la répartition de la masse atmosphérique dictée par la pression.
\end{abstract}

\vspace{0.5cm}
\hrule
\vspace{0.5cm}

\section{Paramètres et Données Climatiques}
Outre les constantes astronomiques standards ($S_0 = 1361 \, \text{W/m}^2$, $\alpha = 0.3$, $\sigma = 5.67 \times 10^{-8} \, \text{W/m}^2/\text{K}^4$), ce modèle intègre des variables chimiques et de pression :
\begin{itemize}
    \item \textbf{Pression de surface ($P_0$)} : $1013.25 \, \text{hPa}$. La pression est utilisée comme proxy de la masse atmosphérique selon l'équation de l'hydrostatique.
    \item \textbf{Concentration en $CO_2$ ($C_{CO_2}$)} : Fixée dans l'exemple à $415 \, \text{ppm}$, représentative des valeurs contemporaines (données NOAA).
    \item \textbf{Coefficient d'absorption ($k_{CO_2}$)} : Paramètre effectif calibré à $0.0035 \, \text{ppm}^{-1}$ pour restituer l'opacité de la bande infrarouge.
\end{itemize}

\section{Discrétisation et Loi de Beer-Lambert}
L'atmosphère est divisée en deux couches de masses strictement égales, c'est-à-dire représentant chacune la moitié de la colonne de pression :
\begin{itemize}
    \item \textbf{Couche 1 (basse)} : de $1013.25 \, \text{hPa}$ à $506.6 \, \text{hPa}$ ($\Delta P_1 = P_0 / 2$).
    \item \textbf{Couche 2 (haute)} : de $506.6 \, \text{hPa}$ à $0 \, \text{hPa}$ ($\Delta P_2 = P_0 / 2$).
\end{itemize}
L'épaisseur optique d'une couche dépend de la quantité de gaz qu'elle contient. L'émissivité $\epsilon_i$ de chaque couche est calculée analytiquement par la loi de Beer-Lambert :
\begin{equation}
    \epsilon_i = 1 - \exp\left(-k_{CO_2} \cdot C_{CO_2} \cdot \frac{\Delta P_i}{P_0}\right)
\end{equation}

\section{Mise en Équation des Bilans Couplés}
Pour simplifier les notations, on pose les flux d'émission du corps noir : $U_0 = \sigma T_0^4$, $U_1 = \sigma T_1^4$ et $U_2 = \sigma T_2^4$. La complexité du modèle à $N$ couches réside dans la \textbf{transmittance} : un flux peut traverser une couche sans être absorbé. La fraction transmise à travers la Couche 1 est $(1 - \epsilon_1)$.

\subsection{Équation 1 : Bilan de la surface terrestre ($T_0$)}
La surface absorbe le flux solaire $F_{sol}$, le rayonnement infrarouge de la Couche 1, \textbf{et} le rayonnement de la Couche 2 qui a réussi à traverser la Couche 1. Elle émet $U_0$ vers le haut.
\begin{equation}
    \text{Bilan}_{surf} = F_{sol} + \epsilon_1 U_1 + \epsilon_2 U_2 (1 - \epsilon_1) - U_0
\end{equation}

\subsection{Équation 2 : Bilan de la Couche 1 ($T_1$)}
La Couche 1 (basse) absorbe une fraction $\epsilon_1$ du rayonnement de la surface (venant du bas) et une fraction $\epsilon_1$ du rayonnement de la Couche 2 (venant du haut). Elle émet son propre rayonnement vers le haut et vers le bas ($2 \epsilon_1 U_1$).
\begin{equation}
    \text{Bilan}_{C1} = \epsilon_1 U_0 + \epsilon_1 (\epsilon_2 U_2) - 2 \epsilon_1 U_1
\end{equation}

\subsection{Équation 3 : Bilan de la Couche 2 ($T_2$)}
La Couche 2 (haute) absorbe le rayonnement de la surface qui a traversé la Couche 1, ainsi que le rayonnement émis par la Couche 1. Elle émet vers l'espace et vers le bas ($2 \epsilon_2 U_2$).
\begin{equation}
    \text{Bilan}_{C2} = \epsilon_2 U_0 (1 - \epsilon_1) + \epsilon_2 (\epsilon_1 U_1) - 2 \epsilon_2 U_2
\end{equation}

\section{Résolution Numérique}
Le système d'équations est résolu par l'algorithme de recherche de racines \texttt{fsolve} de la bibliothèque \texttt{scipy.optimize}. Le vecteur initial (\texttt{T\_guess}) est choisi de manière physiologiquement cohérente avec le gradient thermique troposphérique normal pour assurer une convergence rapide et éviter des minima locaux mathématiques aberrants :
\begin{itemize}
    \item $T_0 \approx 288 \, \text{K}$ (Surface)
    \item $T_1 \approx 260 \, \text{K}$ (Couche basse)
    \item $T_2 \approx 230 \, \text{K}$ (Couche haute)
\end{itemize}

\end{document}